\scanpage{38}
\unsure{Rappelons|Rappel}
\diagram{\[
(42)\qquad\begin{tikzcd}[ampersand replacement=\&,column sep=large]
E(\unsure{\alpha|\omega})\arrow[r,phantom,"\subset"]\arrow[d,phantom,"\Vert" description]\&\pi_\infty\arrow[r,"p_3"]\&\pi_3\\
p_3^{-1}(\alpha)\&\&
\end{tikzcd}
\]}
\[
\alpha=(\overline j,\deleted{\unsure{\overline{\vphantom j}}}j,1)\bmod\mu_3
=\{(\overline j,\deleted{\unsure{\overline{\vphantom j}}}j,1),
(\deleted{\unsure{\overline{\vphantom j}}}j,1,\overline j),
(1,\overline j,\deleted{\unsure{\overline{\vphantom j}}}j)\}
\]
\margin{$\unsure{\in|\cap}\ \pi_3$}
\deleted{$\unsure{j}$}

\deleted{\unsure{Posons}.}
Nous \deleted{prendrons} utiliserons
\[
(\qquad j=\exp\frac{2i\pi}{3},\quad \overline j=\exp-\frac{2i\pi}{3}
\]
et posons
\[
(43)\qquad\zeta_n=\exp\frac{2i\pi}{3n}
\]
et
\[
\deleted{\unsure{\alpha_n}}\ \deleted{\unsure{\zeta_n}}\quad
\alpha_n=(\overline\zeta_n,\deleted{\unsure{\overline{\vphantom\zeta}}}\zeta_n,
1\deleted{\unsure{\zeta_n}})
\]
\deleted{mod $\mu_n$ dans}
\margin{considéré soit comme un point de $\overline C_n$ au dessus de
$\alpha=-\overline j$ soit comme point de $\mu_{3n}^3/\mu_{3n}$
\deleted{\unsure{déf.}} $=\pi_{3n}$}
\[
\alpha_\infty=
\deleted{\unsure{(\zeta_\infty,\overline\zeta_\infty,1)}}
\inserted{(\overline\zeta_\infty,\zeta_\infty,1)}
\]
\margin{considéré soit comme pts de $\overline C_\infty$ au dessus de
$\alpha=-\overline j$ soit comme él. de $\mu_\infty^3/\mu_\infty$
$=\pi_\infty$}
\[
(44)\qquad\zeta_\infty=(\zeta_n)\in T
\]
est la base canonique de $T$, donnée par l'exponentielle complexe,
permettant donc de définir un isom. can.
\diagram{\[
(45)\qquad\begin{tikzcd}[ampersand replacement=\&]
\widehat{\mathbb Z}\arrow[r,"\varphi","\sim"']\&T
\end{tikzcd}
\]}
% archive foliation: 37
